\section{Per-participant cohort and negative construction}
\label{app:cohort}

The modelling set behind every reported result: the ten retained
participants at the fixed negative-to-positive ratio of 1.31 adopted in
Section 4.3, for negative-sampling seed 0. A ratio below the 1.31 target
means the eligible negative pool was exhausted before the target was met;
the lowest, at 1.03, sets the residual balance spread of
$1.28\times$ quoted in Table 3.

\begin{table}[h]
\centering
\begin{tabular}{lrrr}
\toprule
Participant & Positive windows & Negative windows & Ratio \\
\midrule
8B & 74 & 76 & 1.03 \\
5C & 68 & 84 & 1.24 \\
94 & 34 & 44 & 1.29 \\
7A & 248 & 324 & 1.31 \\
6B & 182 & 238 & 1.31 \\
E4 & 310 & 406 & 1.31 \\
F5 & 116 & 152 & 1.31 \\
83 & 314 & 412 & 1.31 \\
15 & 79 & 104 & 1.32 \\
BG & 154 & 203 & 1.32 \\
\midrule
Total & 1579 & 2043 & 1.29 \\
\bottomrule
\end{tabular}
\caption{Cohort and negative construction at the adopted ratio of 1.31, sorted by achieved ratio.}
\end{table}

\section{Exclusion cascade}
\label{app:attrition}

Window counts through the exclusion rules, in the order fixed in advance.
Rule numbering follows the analysis plan; rules that remove no windows at
the window level are omitted.

\begin{table}[h]
\centering
\begin{tabular}{cp{5.8cm}rrr}
\toprule
\# & Rule & Entering & Removed & Surviving \\
\midrule
1 & sessions $<$ 5 min, before non-wear removal & 37252 & 32 & 37220 \\
2 & non-wear windows & 37220 & 87 & 37133 \\
3 & dead-EDA sessions & 37133 & 598 & 36535 \\
8 & participant 6D, no eligible negatives & 36535 & 681 & 35854 \\
9 & flat-EDA participants (four) & 35854 & 11520 & 24334 \\
\bottomrule
\end{tabular}
\caption{Exclusion cascade at the window level.}
\end{table}

\section{Per-fold results}
\label{app:perfold}

Leave-one-participant-out folds, averaged over three negative-sampling
seeds, with the across-seed range. Episode recall is at one false alarm
per worn hour. These are the folds behind the aggregate figures in
Section 5; we report them individually because the spread is wide and the
two lowest-count folds are not comparable to the rest.

\begin{table}[h]
\centering
\begin{tabular}{lrrcr}
\toprule
Participant & Episodes & AUC & AUC range & Episode recall \\
\midrule
5C & 4 & 0.388 & 0.365--0.401 & 0.000 \\
BG & 13 & 0.633 & 0.622--0.646 & 0.000 \\
6B & 10 & 0.682 & 0.671--0.701 & 0.200 \\
15 & 10 & 0.811 & 0.794--0.828 & 0.067 \\
E4 & 23 & 0.813 & 0.804--0.825 & 0.232 \\
94 & 5 & 0.835 & 0.816--0.865 & 0.133 \\
83 & 15 & 0.841 & 0.828--0.861 & 0.089 \\
7A & 23 & 0.847 & 0.827--0.867 & 0.029 \\
F5 & 23 & 0.882 & 0.867--0.898 & 0.159 \\
8B & 12 & 0.885 & 0.858--0.909 & 0.028 \\
\midrule
Mean & 138 & 0.762 & & 0.094 \\
\bottomrule
\end{tabular}
\caption{Per-fold AUC and episode recall, seed-averaged. The fold range quoted in the main text, 0.388 to 0.885, is the range of this column.}
\end{table}

\section{Operating points}
\label{app:operating}

The alarm-rate trade-off, averaged over three seeds. The paper reports the
one-alarm-per-worn-hour row; the others are given so the shape of the
curve is visible rather than asserted.

\begin{table}[h]
\centering
\begin{tabular}{rrrrc}
\toprule
False alarms/h & FPR & Window recall & Episode recall & Range \\
\midrule
0.5 & 0.0176 & 0.087 & 0.029 & 0.029--0.029 \\
1.0 & 0.0357 & 0.155 & 0.106 & 0.058--0.130 \\
2.0 & 0.0715 & 0.292 & 0.280 & 0.261--0.319 \\
\bottomrule
\end{tabular}
\caption{Operating points over 138 episodes with usable coverage.}
\end{table}
